\chapter{Conclusion}

In conclusion, to answer the original research question, it seems that it is indeed possible to deploy a quantized machine learning model on an \gls{esp32} microcontroller to predict \gls{logp} values based on ECFP4 fingerprints with a high degree of accuracy.
The proof-of-concept model achieved an \acrshort{mae} of 0.29, an \acrshort{rmse} of 0.39, and an $R^2$ value of 0.92 on the test set, which indicates that the predictions are closely aligned with the true values.

The inference time of 53.13 milliseconds per compound suggests that the model is efficient enough for real-time applications, although the overall round trip time of 237.97 milliseconds indicates that there may be some overhead in the communication stages that could be optimized in future iterations.

A notable shortcoming of the current model is that the dataset it was trained on was derived from the \gls{ZINC20} database of drug-like compounds, 
and therefore, the predictions may not generalize well to non-drug-like molecules or ones with extreme \gls{lipophilicity} or \gls{hydrophilicity}.
So, a potential future direction for this research would be to train and test a model on a more diverse dataset of compounds with experimentally determined \gls{logp} values.